\begin{styletableII}

Определение
    & \verb|\defi{Метка}{Термин}{Текст}|
        \newline
        \verb|[Примеры][Замечание]|
\\

Утверждение
    & \verb|\stat{Метка}{Название}{Текст}|
        \newline
        \verb|[Док-во][Замечание]|
\\

Лемма
    & \verb|\lemm{Метка}{Название}{Текст}|
        \newline
        \verb|[Док-во][Замечание]|
\\

Теорема
    & \verb|\theo{Метка}{Название}{Текст}|
        \newline
        \verb|[Док-во][Замечание]|
\\

Следствие
    & \verb|\conc{Метка}{Название}{Текст}|
        \newline
        \verb|[Док-во][Замечание]|
\\

Пример
    & \verb|\exam{Метка}{Название}{Условие}|
        \newline
        \verb|[Решение][Замечание]|
\\

Замечание
    & \verb|\note{Метка}{Название}{Текст}|
        \newline
\\

Алгоритм
    & \verb|\algo{Метка}{Название}{Описание}|
        \newline
        \verb|[Сложность][Замечание]|
\\

Код
    & \verb|\code{Метка}{Название}{Файл с кодом}|
        \newline
        \verb|[Файл с выводом]|
\\

Задача
    & \verb|\task{Текст задачи}|
        \newline
\\

\hline
\end{styletableII}

\stylenote{
    Используемый стиль представления определений (\texttt{\textbackslash defi}) предполагает, что в шапке (второй аргумент -- <<Термин>>) записывается непосредственно определяемое понятие, а далее в теле блока уже приводится его определение (третий аргумент -- <<Текст>>), то есть между <<Термин>> и <<Текст>> подразумевается знак тире. 
}

\stylenote{
    Опциональный аргумент <<[Примеры]>> команды \texttt{\textbackslash defi} автоматически оборачивается в окружение \texttt{itemize}, поэтому его содержимое должно состоять из элементов списка, начинающихся с команды \texttt{\textbackslash item}.
}

\stylenote{
    Для блоков допустимы пустые обязательные аргументы (метка, термин или название и текст): если метка или название не нужны, то передаётся пустая пара фигурных скобок.
}

\stylenote{
    Все блоки, кроме \texttt{\textbackslash code} и \texttt{\textbackslash task}, могут использоваться в <<звёздной форме>> (например, \texttt{\textbackslash defi*}), при этом будет создан ненумерованный блок, и счётчик блока не будет увеличиваться.
    На звёздные блоки ссылки не ставятся, то есть первый аргумент всегда должен быть пустым.
}

\stylenote{
    Команда \texttt{\textbackslash code} используется для вставки программного кода из внешнего файла.
    Первый аргумент задаёт метку блока, второй~--- название, третий~--- путь к файлу с исходным кодом.
    В текущей версии оформления блок \texttt{\textbackslash code} предназначен для кода на Python.

    Если необходимо привести не только исходный код, но и вывод программы, путь к файлу с выводом передаётся в необязательном четвёртом аргументе.
    Например:
    \begin{center}
        \small
        \texttt{\textbackslash code\{code:example\}\{Пример\}\{code/l-01/example.py\}[code/l-01/example.txt]}
    \end{center}
    В этом случае вывод программы автоматически печатается внутри того же блока после исходного кода.
    Если необязательный аргумент не указан, раздел с выводом не создаётся.

    Внутри исходного файла Python, передаваемого в \texttt{\textbackslash code}, можно использовать специальную команду \texttt{\textbackslash cmo\{N\}} в комментариях после вызовов \texttt{print}, чтобы связать строку кода с соответствующей строкой вывода программы.
    Команда оформляется как комментарий и заключается в кавычки-ёлочки, например:
    \begin{center}
        \small
        \texttt{print(last\_digit(number))~~~~~\# «\textbackslash cmo\{1\}»}
    \end{center}
    При сборке документа значение \texttt{N} отображается в виде стилизованного бейджа, прижатого к правому краю строки кода.
    Это удобно, когда требуется явно сопоставить строки исходного кода со строками его вывода.
    Сами комментарии вида \texttt{\# «\textbackslash cmo\{...\}»} автоматически удаляются Lua-препроцессором, поэтому в листинге кода остаётся только бейдж.
}